\subsection{Principal norms}\label{sec:norms}
\label{thm:norm-order}\label{thm:norm-leading}

Take \(Q(x,y)=N_{K/\Q}(x+y\omega)\) in a basis \(1,\omega\) of
a fixed order \(\mathcal O\) of discriminant \(\Delta=f^2D\).
Set \(P=P_Q\), \(\Lambda=\Lambda_Q\), \(R_T=R_{T,T}\),
\(S=\{p:p\mid2\Delta\}\), \(\chi=\chi_D\), and \(J_{\mathcal O}=J_S\).
Theorem~\ref{cor:common-field-form-pairs} with both slots equal to
\(Q\) gives the principal order bounds through part~\ref{part:form-classical}
and, under the analytic inputs, the leading asymptotic through
part~\ref{part:form-leading}. Its coefficient is
\begin{equation}\label{eq:norm-coefficient}
 \mathfrak A_{\mathcal O}(N)=
 \frac{4L(1,\chi)}{\pi}
 \prod_{p\in S}\frac{1-\chi(p)/p}{1-1/p}
 \prod_{\substack{p\notin S\\\chi(p)=-1}}
                  \left(1-\frac2{p(p-1)}\right)
 \prod_{p\in S}d_p(N).
\end{equation}
Here \(d_p=d_{p;1,2}\) for \(Q_1=Q_2=Q\).
The norm form has \(A=\mathcal O\), hence \(c_Q=0\), so
Lemma~\ref{lem:form-ideal-shadow} becomes
\begin{equation}\label{eq:norm-shadow-ideals}
 n\in\Lambda\quad\Longleftrightarrow\quad
 \text{some primitive integral invertible }I
       \text{ has }n(I)=n,\ [I]\in2\operatorname{Pic}(\mathcal O);
\end{equation}
the condition \([I]=0\) characterizes \(P\). In particular,
Proposition~\ref{prop:norm-class-error} gives
\(R_\Lambda-R_P\ll_{\mathcal O}NJ_{\mathcal O}/(\log N)^{6/5}\),
and \(P=\Lambda\) when the doubled class group vanishes.

\begin{remark}[Gaussian normalization]
For $\mathcal O=\Z[i]$ one has $P=C_4\sqcup2C_4$ and
\[
 d_2(N)=
 \begin{cases}
 1/4,&N\bmod8\in\{2,3,6,7\},\\
 1/8,&N\bmod8=4,\\
 0,&N\bmod8\in\{0,1,5\}.
 \end{cases}
\]
Consequently the coefficient of $J_{\mathcal O}(N)N/\log N$ is
$\frac12\prod_{p\equiv3\ (4)}(1-2/(p(p-1)))$ in the first four
classes and half that coefficient in the class $4\bmod8$.
This also calibrates the value-density normalization: the two Gamma
factors are already included in $c_{\mathcal C}^2$, and no additional
beta-integral factor occurs.
\end{remark}
